\section{Contesto: V2X e guida cooperativa}

La cooperazione tra veicoli rappresenta una delle direttrici più rilevanti nello sviluppo degli Intelligent Transportation Systems, perché consente di superare i limiti della percezione locale del singolo mezzo e di coordinare in modo più stretto decisioni di controllo, sicurezza ed efficienza. In questo quadro, le comunicazioni Vehicle-to-Everything (V2X) hanno un ruolo dominante e diventano parte integrante della catena di controllo, poiché il valore informativo di un messaggio dipende dal fatto che esso arrivi con tempi e affidabilità compatibili con la dinamica del veicolo e con i margini di sicurezza richiesti dall’applicazione \cite{dressler2019cooperative,giordano2019joint,segata2016safe}.

Tra le applicazioni cooperative più esigenti rientra la guida in \emph{platooning}, in cui più veicoli viaggiano coordinati a distanza ridotta scambiando periodicamente informazioni sul proprio stato. In questo scenario emerge con particolare evidenza il limite delle soluzioni mono-tecnologia: congestione, interferenza, occlusioni ottiche, \emph{blind spot} e handover possono degradare il link e propagare l’effetto del degrado fino al controllore longitudinale, con impatti diretti su comfort, prestazioni e sicurezza \cite{segata2015communication,bazzi2020wireless,segata2023multi-technology}.

Questa tesi assume come riferimento sperimentale il lavoro di Segata \emph{et al.} \cite{segata2023multi-technology}, che propone un’architettura multi-tecnologia e il meccanismo \emph{SafeSwitch} per gestire in modo sicuro i passaggi tra diversi livelli di affidabilità di rete. Il contributo del presente elaborato consiste nel riprodurre e aggiornare tale impostazione su una toolchain più recente, basata su versioni aggiornate di OMNeT++, Veins e PLEXE, e nell’integrare Simu5G al posto della precedente componente cellulare basata su SimuLTE, così da verificare la trasferibilità dell’approccio in un ambiente coerente con l’evoluzione attuale degli strumenti di simulazione \cite{nardini2020simu5g,virdis2014simulte}.

Tale scelta assume particolare rilevanza alla luce dell’evoluzione delle tecnologie V2X. Le più recenti specifiche di C-V2X prevedono infatti un progressivo passaggio dalle implementazioni iniziali basate su LTE verso soluzioni fondate su reti 5G. L’integrazione di Simu5G consente quindi di valutare l’architettura proposta in un contesto più aderente agli sviluppi attuali delle comunicazioni veicolari.
